% Einheit 2 von 2 – Einheit 2 (bank/vierfeldertafel/e2.jsonl, zusammenbau v0.1)
%% TODO Prüfungswort Einheit 2: nicht in der Marken-Zeile
%% TODO Zeitmarke Einheit 2: Marken-Zeile nicht lesbar
%% TODO Zweigzeile Teil 1: was hier gelernt wird (Satz in Schülersprache aus dem Einheitstitel) – Einheit 2
%% TODO Einheitentitel 2 nicht in der Mappe lesbar
\einheitenkopf[e2]{Einheit 2 von 2 · Einheit 2}
\setcounter{aufgabe}{10}

%% TODO Titel: Ich-kann-Satz fehlt in der Bank (Platzhalter: Kettenname) – Nr. 11
%% TODO Anweisung: ein Satz über der ersten Teilaufgabe fehlt in der Bank – Nr. 11
\begin{aufgabe}{Aus der Tafel rechnen}
%% TODO \rechenplatz{n}: Schrittzahl des Grundfalls fehlt in der Bank (nur bei fester Schreibform, 2.3 b)
\begin{teile}
\teil $A$: Ein Kunde kauft Brot, $B$: er kauft Brötchen. Gesucht ist die Wahrscheinlichkeit für „Brot oder Brötchen“. Welcher Ausdruck passt? \\ \kreuz{$1 - P(\overline{A} \cap \overline{B})$} \\ \kreuz{$P(A \cap \overline{B}) + P(\overline{A} \cap B)$} \\ \kreuz{$P(\overline{A} \cap \overline{B})$}
\teil Von $300$ Gästen eines Hotels sind $120$ Familien ($F$); $180$ Gäste buchen Halbpension ($H$), davon $90$ Familien. Wie viele Gäste sind keine Familie und buchen keine Halbpension? \leerfeld
\teil Von den Studierenden einer Hochschule wohnen $58\,\%$ in der Stadt ($W$), $30\,\%$ studieren ein technisches Fach ($T$), und $21\,\%$ tun beides. Mit welcher Wahrscheinlichkeit wohnt eine zufällig ausgewählte Person in der Stadt oder studiert ein technisches Fach? \hfill (Abitur 2023 GK) \\ \leerfeld[\%]
\teil Bei einem Sportfest sind $70\,\%$ der Teilnehmer Kinder, die übrigen Erwachsene. $14\,\%$ aller Teilnehmer sind Erwachsene mit Medaille; von den Kindern hat die Hälfte eine Medaille. Welcher Anteil aller Teilnehmer hat keine Medaille? \hfill (Abitur 2020 GK) \\ \leerfeld[\%]
\teil Die Vierfeldertafel zeigt für die Bevölkerung einer Großstadt die Anteile in Prozent ($S$: Die Person hat ein Streaming-Abo, $H$: sie wohnt in einem Haus mit Garten). Aussage: „$P(S \cap \overline{H})$ ist etwa halb so groß wie die Wahrscheinlichkeit, dass entweder $S$ oder $\overline{H}$ eintritt.“ Beurteile die Aussage. \hfill (Abitur 2025 LK)

\vierfeldertafel{H}{S}{18,27,45,22,33,55,40,60,100}
\end{teile}
\end{aufgabe}

%% TODO Titel: Ich-kann-Satz fehlt in der Bank (Platzhalter: Kettenname) – Nr. 12
%% TODO Anweisung: ein Satz über der ersten Teilaufgabe fehlt in der Bank – Nr. 12
\begin{aufgabe}{Aus der Tafel rechnen – Fehler finden, Begründen}
\begin{teile}
\teil Ich finde den Fehler: $50\,\%$ der Mitglieder einer Bücherei lesen Krimis ($K$), $40\,\%$ Sachbücher ($S$), $15\,\%$ beides. Ida rechnet: \rechnung{P(K \cup S) &= 50\,\% + 40\,\% = 90\,\%}
\teil Begründe, warum „$A$ oder $B$“ das Gegenereignis von „weder $A$ noch $B$“ ist.
\end{teile}
\end{aufgabe}
